% SPECTRA manuscript. Build with `make` (latexmk or tectonic).
\documentclass[11pt]{article}

% ---------- layout and fonts (readability first) ----------
\usepackage[letterpaper,margin=0.75in]{geometry}   % 7.0 in text width = figure width
\usepackage[T1]{fontenc}
\usepackage{lmodern}
\usepackage{microtype}
\usepackage{setspace}
\setstretch{1.2}
\usepackage[parfill]{parskip}
\usepackage{amsmath}
\usepackage{textcomp}
\usepackage{graphicx}
\usepackage{float}
\usepackage{booktabs}
\usepackage{tabularx}
\usepackage{array}
\newcolumntype{L}[1]{>{\raggedright\arraybackslash}p{#1}}
\newcolumntype{C}[1]{>{\centering\arraybackslash}p{#1}}
\usepackage[labelfont=bf,labelsep=period,font=small,justification=justified,singlelinecheck=false]{caption}
\usepackage{authblk}
\renewcommand\Authfont{\normalsize}
\renewcommand\Affilfont{\small}
\usepackage[super,sort&compress,comma]{natbib}
\usepackage[hidelinks]{hyperref}
\usepackage{lineno}
\linenumbers                                       % comment out for a clean reading copy

% Figures: public repository layout first, then project layout.
\graphicspath{{../figures/main/}{../figures/supplementary/}{../figures/manuscript_20261111/final_figures/}}

\newcommand{\fig}[1]{\includegraphics[width=\textwidth,height=0.72\textheight,keepaspectratio]{#1}}
\newcommand{\pnl}[1]{\textbf{#1}}                  % panel letter in captions

\title{\textbf{Attention-based learning in peripheral blood reveals cancer signals beyond circulating tumor cells}}

\author[1,2]{Dean Tessone}
\author[1,3]{Alireza Rezvani}
\author[1,2]{Ahmed Saadawy}
\author[1,2]{Niki Hakami}
\author[1,2]{Isabel Romero}
\author[1,2]{Doanna Pham}
\author[1,2]{Mina Tolou Ghani}
\author[1,2]{Ian Ward}
\author[1]{Elizabeth Qi}
\author[1]{Andres Rivera}
\author[1]{Avery Lin}
\author[1]{Rafael Nevarez}
\author[1]{Will Kim}
\author[1]{Valerie Hennes}
\author[1,4]{Jorge Nieva}
\author[1]{Stephanie Shishido}
\author[1]{Jeremy Mason}
\author[1,5]{Assad A Oberai}
\author[1,2,4]{James Hicks}
\author[1,2,3,4,5,6]{Peter Kuhn}
\affil[1]{Convergent Science Institute for Cancer, University of Southern California, Los Angeles, CA}
\affil[2]{Department of Molecular and Cellular Biosciences, Dornsife College of Letters, Arts, and Sciences, University of Southern California, Los Angeles, CA}
\affil[3]{Department of Biomedical Engineering, Viterbi School of Engineering, University of Southern California, Los Angeles, CA}
\affil[4]{Norris Comprehensive Cancer Center, Keck School of Medicine, University of Southern California, Los Angeles, CA}
\affil[5]{Department of Aerospace and Mechanical Engineering, Viterbi School of Engineering, University of Southern California, Los Angeles, CA}
\affil[6]{Catherine and Joseph Aresty Department of Urology, Keck School of Medicine, University of Southern California, Los Angeles, CA}
\date{}

\begin{document}
\maketitle

\begin{abstract}
\noindent Blood carries cells shed by tumors and by the host's response to them, capturing information about cancer and its metastatic cascade. Unlike cell-free DNA, intact circulating cells preserve the morphology, proteins and genomes of single cells, offering direct access to the cellular mechanisms of metastasis. Yet enrichment-based workflows capture one predefined cell type at a time, typically cytokeratin-positive circulating tumor cells. Here we present SPECTRA, a weakly supervised transformer that learns from enrichment-free images of rare circulating cells using only patient-level diagnoses, and traces each prediction back to individual cells. Across 13 cancer types, SPECTRA separated 713 held-out cancer patients from 44 non-cancer controls (AUROC 0.82), whereas counts of cytokeratin-positive candidate cells performed at chance (AUROC 0.50). Attribution identified 23 recurrent cell populations, led by large, vimentin-high cells. Among a selected subset of these cells profiled by imaging mass cytometry and single-cell sequencing, most were copy number neutral and carried endothelial or myeloid markers. Together, enrichment-free, weakly supervised imaging thus detects cancer-associated cellular information beyond canonical cytokeratin-positive circulating tumor cells, integrating contributions from multiple cellular populations.
\end{abstract}
\clearpage

% =====================================================================
\section*{Introduction}

Advances in liquid biopsy have established peripheral blood as a minimally invasive source of clinically actionable cancer measurements. As of the time of writing, most translational progress has centered on circulating tumor DNA (ctDNA) mutational profiling, methylation signatures, or fragmentation\cite{bettegowda2014,liu2020,cristiano2019}. While ctDNA-based approaches have achieved substantial success in cancer detection and monitoring, they do not preserve intact tumor states, clonal heterogeneity, protein expression, or the phenotypic diversity of tumor and microenvironmental cells in circulation, which are necessary for better understanding the circulatory cancer biology. Interrogating the cellular compartment of blood offers a complementary dimension to liquid biopsy uniquely suited to resolving single-cell biology.

The cellular compartment of blood enables access to intact biological material, linking high-resolution morphological and protein-level phenotypes to downstream single-cell genomic and multiplexed proteomic profiling\cite{dago2014,gerdtsson2018}. Beyond circulating tumor cells (CTCs), which capture tumor heterogeneity and epithelial-to-mesenchymal transition\cite{yu2013}, peripheral blood contains diverse non-malignant and hybrid populations associated with the tumor microenvironment, including circulating endothelial cells\cite{bertolini2006}, tumor-associated macrophage-like cells\cite{adams2014}, cancer-associated fibroblasts\cite{ao2015}, and cells co-expressing epithelial and hematopoietic markers\cite{gast2018,higa2025,naghdloo2025caught}.

Because conventional enrichment workflows isolate cells based on predefined surface markers or biophysical properties\cite{allard2004,nagrath2007,vona2000,ozkumur2013,sollier2014}, they inherently discard subsets outside those criteria and make it impossible to survey all relevant populations simultaneously. High-throughput, whole-blood imaging enables these rare populations to be surveyed without prior enrichment\cite{marrinucci2012}, yet identifying candidate cells remains a major bottleneck\cite{murgoitio2025,naghdloo2025blue}. Morphologically or phenotypically atypical cells occur in both cancer patients and non-cancer controls, and establishing significance typically depends on marker-gating (CK+, CD45$-$), expert visual curation, post-hoc association with clinical variables, and orthogonal molecular validation\cite{bai2024,resnick2024,malihi2020,shishido2025}. Relying exclusively on \textit{a priori} defined criteria of `interesting' biology risks overlooking biologically informative populations whose phenotypes fall outside established conventions.

Addressing this bottleneck requires shifting from rule-based cell selection to data-driven discovery guided directly by patient-level clinical status. Weakly supervised deep learning has transformed computational pathology and related high-content imaging fields by solving an analogous problem: enabling models to learn diagnostic features from sample-level labels without requiring exhaustive manual annotation of individual cells or tissue regions\cite{campanella2019,lu2021}. Across gigapixel whole-slide images, formulating tissues as unordered collections of image patches has allowed weakly supervised frameworks not only to achieve clinical-grade diagnostic performance, but also to localize tumor foci and uncover previously unrecognized morphological correlates of disease progression and treatment response\cite{campanella2019,lu2021,courtiol2019,skrede2020}. Transformer self-attention allows each cell's representation to be updated in the context of the other cells within a sample\cite{vaswani2017,lee2019,shao2021,wagner2023}. Crucially, the value of this formulation in our setting extends beyond sample-level classification: attention and attribution scores trace cohort-level predictions back to the specific instances driving each decision, providing a principled way to prioritize rare circulating cells for experimental recovery and molecular validation without defining every relevant phenotype \textit{a priori}.

Here, we present SPECTRA, a weakly supervised Transformer framework that connects cohort-scale learning across peripheral blood images with the discovery and molecular characterization of cancer-associated circulating cells. Trained solely on patient-level diagnosis across 1,793 individuals spanning 13 cancer types, SPECTRA distinguishes held-out cancer patients from non-cancer controls (AUROC 0.82) and uses attention attribution to trace each sample-level prediction back to the individual cells carrying the disease signal. Across the cohort, these high-attribution cells organize into 23 recurrent populations. Although classical cytokeratin-positive, CD45$-$/CD31$-$ circulating tumor cells appear among these subsets and are enriched in cancer, their counts alone do not separate patients, and removing every CK-gated cell leaves model performance unchanged. Instead, SPECTRA integrates contributions from a broader cellular repertoire, most prominently large vimentin-high cells alongside aneuploid cytokeratin-positive and dual-positive cells. Orthogonal profiling of a selected subset via imaging mass cytometry and single-cell sequencing shows that, among profiled cells, vimentin-high populations are mostly copy-number neutral and include CD31+ endothelial-like and myeloid-marker-positive cells. These populations emerge directly from patient-level supervision without predefined gating rules, and their relative weighting shifts when the model is adapted to more challenging clinical comparisons. By linking cohort-scale learning to targeted single-cell recovery, SPECTRA provides a framework for discovering cancer-associated cellular phenotypes beyond conventional CTC definitions, adding a biologically interpretable cellular layer to liquid biopsy.

% =====================================================================
\section*{Results}

\subsection*{The SPECTRA framework and multi-cancer cohort}

To identify cancer-associated circulating cells without predefined gating criteria, we developed SPECTRA, an end-to-end framework that combines enrichment-free, fluorescent whole-blood imaging with attention-based set deep-learning models (Fig.~\ref{fig:1}a,b). Following red blood cell lysis, nucleated cells from peripheral blood (approximately 5 million cells, corresponding to a mean blood volume of 1.5~mL) are deposited onto glass slides without marker or size-based enrichment. Cells are subsequently immunofluorescently stained for DAPI (DNA), pan-cytokeratin (CK, epithelial), vimentin (Vim, mesenchymal), and CD45/CD31 (leukocyte and endothelial, respectively), and imaged via automated whole-slide scanning (Fig.~\ref{fig:1}a). Individual cells are segmented and encoded into 512-dimensional feature vectors using a pre-trained, frozen self-supervised encoder, trained on 3.8 million rare-event images without diagnosis labels (Supplementary Fig.~\ref{sfig:1}a)\cite{naghdloo2025blue}. Unsupervised anomaly detection isolates the most phenotypically atypical group of cells in each sample ($<$0.5\% of all nucleated events)\cite{murgoitio2025}. Automated quality control models remove blurry cells and staining artifacts, refining this pool to a median of 1,426 candidate cells per patient (range 21 to 9,129; Supplementary Fig.~\ref{sfig:1}b). SPECTRA represents these embedded cells as an unordered set of tokens, prepended with a learned summary token (Fig.~\ref{fig:1}b). A six-layer Transformer is trained to predict the patient-level diagnosis (cancer vs no-cancer) from the summary token, such that single cells and the summary token attend bidirectionally to pool tumor-defining signals. Attribution modules can subsequently decompose each sample-level prediction back onto individual cells to yield single-cell attribution scores (Fig.~\ref{fig:1}b).

\begin{figure}[!tbp]
\centering
\fig{Figure1.pdf}
\caption{\textbf{Overview of the SPECTRA framework and dataset.} \pnl{a}, Sample preparation and image-processing pipeline. Nucleated cells are isolated from peripheral blood following red blood cell lysis, plated onto glass slides ($\sim$2.5 million cells per slide), stained with a four-channel immunofluorescence assay (DAPI, pan-cytokeratins (CK), vimentin, and CD45/CD31), and scanned to generate a fluorescence whole-slide image (fWSI). Two slides are assessed per patient, for a total of $\sim$5 million analyzed cells (mean blood volume of 1.5~mL). Individual cells are segmented and encoded into feature vectors using a pre-trained deep learning model, followed by unsupervised anomaly detection and quality filtering to isolate rare candidate events ($<$0.5\% of nucleated cells) from each patient sample. \pnl{b}, SPECTRA model overview. SPECTRA is a six-layer, pre-norm Transformer with four attention heads per layer. A learnable patient-level summary token is prepended to the unordered set of rare single-cell embeddings, allowing all cells and the summary token to attend bidirectionally to one another to generate a patient summary representation and predicted cancer probability (SPECTRA score). Attribution traces each sample-level prediction back to individual cells. \pnl{c}, Dataset distribution across 13 cancer types and non-cancer controls ($n=2{,}550$ individuals), stratified into training ($n=1{,}793$) and held-out test ($n=757$) sets. UTUC, upper tract urothelial carcinoma.}
\label{fig:1}
\end{figure}

We applied this pipeline to an archival multi-institutional cohort of 2,550 individuals, comprising 2,402 patients with cancer across 13 cancer types and 148 non-cancer controls (106 healthy donors, 21 women sampled after a cancer-negative mammogram, and 21 individuals with benign pancreatic disease; Fig.~\ref{fig:1}c and Table~\ref{tab:1}). Patients were stratified by cancer type and diagnosis into training and test sets of 1,793 individuals and 757 individuals, respectively, to ensure proper representation of cancer types in both training and evaluation.

\subsection*{SPECTRA detects cancer in peripheral blood}

We evaluated the trained SPECTRA model on our held-out test cohort of 757 individuals (713 patients across 13 cancer types and 44 non-cancer controls). Across all cancers, SPECTRA separated cancer patients from non-cancer controls with an area under the receiver operating characteristic curve (AUROC) of 0.82 [95\% CI: 0.75--0.89] (Fig.~\ref{fig:2}a). SPECTRA maintained consistent discrimination across distinct solid and hematologic malignancies in the test set (Fig.~\ref{fig:2}b). Among major cancer types, SPECTRA performed strongest in lung cancer, achieving an AUROC of 0.90 ($n=62$) (Fig.~\ref{fig:2}b). SPECTRA additionally performed strongly in pancreatic cancer (AUROC 0.88, $n=19$), breast cancer (AUROC 0.84, $n=218$), upper tract urothelial carcinoma (AUROC 0.83, $n=44$), bladder cancer (AUROC 0.82, $n=52$), prostate cancer (AUROC 0.81, $n=264$), and multiple myeloma (AUROC 0.76, $n=23$) (Fig.~\ref{fig:2}b). In matched three-fold cross-validation on the training split, SPECTRA outperformed five other multiple-instance learning approaches, although its margin over the strongest baselines was small (mean fold AUROC 0.81 versus 0.75--0.79; Supplementary Fig.~\ref{sfig:2}a,b).

\begin{figure}[!tbp]
\centering
\fig{Figure2.pdf}
\caption{\textbf{SPECTRA cancer classification performance.} \pnl{a}, Area under the receiver operating characteristic curve (AUROC) for all cancers (AUROC 0.82, 95\% CI 0.75--0.89). \pnl{b}, Cancer-type stratified AUROC curves for lung, pancreatic, breast, upper tract urothelial carcinoma (UTUC), bladder, prostate, and multiple myeloma (and related plasma cell malignancies). \pnl{c}, Patient-level log-odds scores output by SPECTRA. Two thresholds for specificity (90\%, 40 of 44 controls below, and 95\%, 42 of 44) are shown, with patients surpassing the 90\% threshold indicated. Median scores for non-cancer controls clustered around 0, whereas several carcinomas display median scores above the 90\% specificity threshold. UTUC, upper tract urothelial carcinoma. $n$ displayed per row.}
\label{fig:2}
\end{figure}

Examining patient-level SPECTRA scores (log-odds) across the entire cohort showed that non-cancer controls clustered around a median score of $-0.10$ (probability of 0.47), whereas the median score for cancer is 5.85 (probability 0.997) and most cancer types exceeded the 90\% specificity threshold (equivalent to 40 of 44 controls below the threshold, Fig.~\ref{fig:2}c). Patient scores were unrelated to the number of candidate cells (Supplementary Fig.~\ref{sfig:2}d). At 90\% specificity, SPECTRA identified 56.5\% of test cancer patients (403 of 713, 95\% CI 52.9--60.1\%), and 39.7\% at 95\% specificity (283 of 713; 95\% CI 36.2--43.3\%) (Fig.~\ref{fig:2}c).

Because most cancer and control samples came from different clinical studies, we tested whether scores tracked study, collection site, or image acquisition rather than disease (Supplementary Fig.~\ref{sfig:sites} and Supplementary Table~\ref{stab:2}). Discrimination held against each of the four control sources individually (AUROC 0.77--0.88), and excluding any single control source changed pooled AUROC by at most 0.03. Within a single diagnostic work-up study from the same site and medical team, patients with pancreatic cancer were distinguished from individuals with benign pancreatic disease (AUROC 0.89, 95\% CI 0.70--1.00; 0.90 in training CV). Among patients with the same cancer type, we did not detect site-associated differences in scores for hormone-receptor-positive breast cancer (seven sites; site medians 4.5--6.4 log-odds, Kruskal--Wallis $P=0.28$, $\varepsilon^2=0.01$) or lung cancer (five sites; medians 6.1--6.7, $P=0.55$, $\varepsilon^2<0.01$) (Supplementary Fig.~\ref{sfig:sites}c). Model performance was unchanged when restricted to cancer--control pairs matched on scanner instrument, scan year, scan date, or exposure times, and the score's association with cancer was unchanged after adjustment for these variables (Supplementary Fig.~\ref{sfig:2}f). Together, these analyses support that measured pre-analytical variables do not explain model performance.

\subsection*{Attention attribution uncovers 23 recurrent circulating cell phenotypes}

To determine which cellular phenotypes drive SPECTRA predictions, we applied layer-wise relevance propagation (AttnLRP) across all 757 held-out test patients\cite{bach2015,achtibat2024}. For each patient, we isolated the minimal set of cells accounting for 90\% of positive prediction relevance (mass-90), yielding 204,131 high-attribution cells across the test cohort (mean: 269.7 cells/patient). Attribution was highly concentrated, with a median of 15 cells (1\% of the input) expressing half of positive evidence in cancer patients (Supplementary Fig.~\ref{sfig:3}). Leiden clustering of these cells identified 23 recurrent cellular populations spanning diverse morphological and marker profiles (Fig.~\ref{fig:3}a,b)\cite{traag2019}. We annotate each cluster by median fluorescence intensity and cell area, revealing distinct groups encompassing classical cytokeratin-positive, CD45/CD31-negative (CK+, CD45/31$-$) CTC-like events (clusters 20 and 21), dual-positive CK+, CD45/31+ cells (clusters 17, 18, and 19), large vimentin-high (Vim+) cells with variable CD45/CD31 expression (cluster 10), large CD45/CD31-high cells (clusters 12, 13, and 15), among other cell types (Fig.~\ref{fig:3}b). Additional examples of attributed cells from each cluster are shown in Supplementary Fig.~\ref{sfig:4}.

\begin{figure}[!tbp]
\centering
\fig{Figure3.pdf}
\caption{\textbf{Biological and phenotypic characterization of high-attribution circulating cells.} \pnl{a}, Uniform manifold approximation and projection (UMAP)\cite{mcinnes2018} of all mass-90 high-attribution cells across the held-out test cohort ($n=204{,}131$ cells from 757 individuals), colored by 23 Leiden clusters. Representative four-channel microscopy images for each of the 23 clusters (DAPI in blue, pan-cytokeratin (CK) in red, vimentin in white, and CD45/CD31 in green). Images captured with a 40x objective unless indicated as 10x. Scale bars are 5~$\mu$m. \pnl{b}, Cluster-specific marker profiles and attribution metrics. Left: median marker profiles for each cluster with a short name based on its marker profile. Center: AUROC of per-patient cluster abundance distinguishing cancer patients ($n=713$) from controls ($n=44$), where teal bars denote cancer enrichment (AUROC $>$0.5) and orange bars denote control enrichment ($<$0.5). * indicates statistical significance (adjusted $P<0.05$) in two-sided Mann--Whitney U test. Right: model emphasis, $\log_2$ of a cluster's mean per-patient share among attributed cells relative to its mean per-patient share among all input cells, in cancer patients. Blue bars indicate up-weighting by SPECTRA whereas gray bars indicate down-weighting. \pnl{c}, Comparison of an L2-regularized logistic regression on cluster composition alone (AUROC 0.75) against SPECTRA (0.82). Statistical significance assessed by paired bootstrap analysis (4,000 replicates) with $P=0.03$.}
\label{fig:3}
\end{figure}

We next asked whether these recurrent clusters distinguish cancer from non-cancer patients by their abundance alone. Most populations were present at similar frequencies in both groups (18/23 populations with AUROCs of 0.40--0.60). Three clusters showed significant differences ($q<0.05$): small, DAPI-bright events (cluster 1, AUROC 0.69) and CK+, CD45/CD31$-$ CTC-like cells (cluster 20, AUROC 0.64) were more common in cancer patients than controls, while small leukocyte-like CD45/CD31+ cells (cluster 5) were enriched in controls. A classifier trained on each patient's cluster composition alone, using the same 1,793 training patients as SPECTRA, reached a test AUROC of 0.75 (95\% CI 0.67--0.83), below the 0.82 achieved by SPECTRA (paired difference 0.07, 95\% CI 0.01--0.15; $P=0.03$; Fig.~\ref{fig:3}c).

We then asked which populations SPECTRA relies on beyond their frequency. Model emphasis was only weakly correlated with abundance differences (Spearman $\rho=0.38$, $P=0.07$ across clusters). Interestingly, large, vimentin-high cells with variable CD45/CD31 (cluster 10) were the most strongly up-weighted population ($\log_2$ emphasis $+0.88$, 17.6\% attributed relevance (15\%--19\% across model seeds), with stable cluster recovery across Leiden seeds (Jaccard $>$0.89)), despite their abundance showing no significant difference between cancer patients and controls (AUROC 0.47). Dual-positive CK+, CD45/31+ cells (clusters 18 and 19, $\log_2$ emphasis $+0.66$ and $+0.76$) were also up-weighted, with cluster 18 even less frequent in cancer patients (AUROC 0.40). CTC-like clusters 20 and 21 were both enriched and moderately up-weighted ($+0.40$ and $+0.36$). Conversely, the abundant CD45/31-dim (cluster 5; $-1.30$) and CD45/31+ (cluster 13; $-1.05$) populations were strongly down-weighted. Emphasis was largely preserved in non-cancer controls (Spearman $\rho=0.80$ between cancer and control emphasis). This suggests that SPECTRA examines the same cell populations in every patient, and that its prediction depends on the features of those cells rather than on their presence alone.

Consistent with this, cells from the same population had different embeddings in cancer patients and controls. Classifying patients by only the average embedding of their cells within a single cluster, which removes any effect of cell number, gave AUROCs above 0.5 in every cluster (0.53--0.73), significantly in 10 of 23 ($P<0.05$; 8 after correction for multiple testing; Supplementary Fig.~\ref{sfig:5}), including clusters whose abundance did not differ, such as cluster 2 (abundance AUROC 0.53, embedding AUROC 0.73). For cluster 10, the most heavily weighted population, the difference was not captured by measurable cell features (AUROC 0.57 with 74 CellProfiler-style measurements of shape, intensity, or texture) while embeddings showed slightly improved separation (AUROC 0.61). These findings are consistent with SPECTRA drawing on variation within populations as well as their abundance.

\subsection*{\raggedright SPECTRA detects cancer-associated information beyond canonical CK-positive CTCs}

SPECTRA places strong emphasis on circulating cell populations spanning both canonical epithelial populations (clusters 18 and 19) and non-epithelial atypical cells (cluster 10). To test whether SPECTRA relies on classical epithelial CTCs, we compared sample-level predictions with counts of putative CTCs among the same candidate cells that SPECTRA reads (Fig.~\ref{fig:4}). We leveraged two methods to count CTCs in patients: marker-based gating (CK+, CD45/31$-$)\cite{allard2004} and pre-trained cell-type classifiers for epithelial CTCs (CTCs only), immune-like CTCs, and platelet-coated CTCs (all three populations being CTC-like)\cite{naghdloo2025blue}. Of a median of 1,426 cells read by SPECTRA per patient, a median of two were CK-gated (at least one in 32 of 44 controls), while about 600 were classified as CTC-like (Supplementary Fig.~\ref{sfig:1}c). While SPECTRA separates cancer patients from controls with an AUROC of 0.82, enumeration of CTCs by gating scored no better than chance (AUROC of 0.50, and 0.47--0.56 across gating thresholds), with similar performance for putative CTCs only (0.47) and CTC-like classes (0.49) (Fig.~\ref{fig:4}a). Stratifying test patients by the number of CK-gated cells per patient showed that 197 of 713 cancer patients (27.6\%) had zero CK-gated cells detected on their slides (Fig.~\ref{fig:4}b). Among cancers, scores rise slightly with CK-gated count: median 5.4 log-odds with none versus 6.3 with more than ten (Kruskal--Wallis $P=0.009$; Fig.~\ref{fig:4}b). The 197 with none still separate from controls (AUROC 0.80).

Consistent with this independence from classical CTC abundance, SPECTRA maintained strong performance across clinical stages in held-out breast cancer patients with available staging annotations (Fig.~\ref{fig:4}c). Canonical CK-gated (AUROC 0.47 in early stage vs.\ 0.53 in late stage) and CTC-like counts (AUROC 0.52 vs.\ 0.57) failed to distinguish early-stage breast cancer ($n=50$) from controls ($n=44$) and remained near chance in late-stage metastatic disease ($n=98$) (Fig.~\ref{fig:4}c). By contrast, SPECTRA achieved an AUROC of 0.82 (95\% CI 0.73--0.90) in early-stage breast cancer and 0.86 (0.78--0.93) in late-stage breast cancer, demonstrating sensitivity to circulating disease signals prior to widespread metastasis (Fig.~\ref{fig:4}c).

To directly test dependence on specific cell subsets in SPECTRA's predictions, we performed an \textit{in silico} ablation experiment across 757 test individuals and compared each subset removal with removing the same number of randomly chosen cells (Fig.~\ref{fig:4}d and Supplementary Table~\ref{stab:3}). Removing every CK-gated cell, or the CK+, CD45/31$-$ clusters 20 and 21, left discrimination unchanged (AUROC 0.82). Removing all CTC-like cells lowered AUROC to 0.73 (0.80 after matched random removal), mainly by raising scores in controls, indicating that SPECTRA uses information associated with these broader epithelial and hybrid phenotypes (Fig.~\ref{fig:4}d). In contrast, ablating just the top 10 highest attribution cells per patient reduced sensitivity at the original 90\% specificity threshold to 29\% and removing the top 50 attributed cells reduced sensitivity to 15\%, whereas removing an equal number of randomly chosen cells had no effect (Fig.~\ref{fig:4}d). AUROC, however, was preserved (0.84 for top-10, 0.82 for top-50, and 0.83 for cluster 10). The most highly attributed cells raise the overall level of the score, but do not impact patient rankings in our small test population. Assessing the cells with the highest attribution in four patients with no CK-gated cells revealed a distinct population of large, vimentin-positive cells, with variable CD45/CD31, all belonging to cluster 10 (Fig.~\ref{fig:4}e). Together, these results show that SPECTRA's cancer signal integrates information from many cell populations, including putative CTCs and novel atypical populations.

\begin{figure}[!tbp]
\centering
\fig{Figure4.pdf}
\caption{\textbf{SPECTRA detects cancer signals beyond canonical circulating tumor cells.} \pnl{a}, Held-out test ROC curves comparing SPECTRA (AUROC 0.82) against patient-level counts of CK-gated cells (0.50), CTC-like cells (0.49), and canonical putative CTCs only (0.47) across $n=713$ cancer patients and $n=44$ controls. \pnl{b}, Violin plots of SPECTRA patient log-odds scores stratified by the number of CK-gated cells per patient (0, 1--2, 3--10, and $>$10). \pnl{c}, Classification AUROC in held-out breast cancer patients stratified by clinical stage (early-stage $n=50$, late-stage $n=98$) for SPECTRA, CK-gated cells, and CTC-like cell counts. \pnl{d}, \textit{In silico} ablations showing held-out test sensitivity at 90\% specificity after removing all CK-gated cells, all CTC-like populations, the top 10 attributed cells, or the top 50 attributed cells (teal), as compared to an equivalent number of randomly selected cells (gray), and the original model baseline (dashed line). Sensitivity is shown at the original locked cut-off because removing top-attributed cells lowers scores in controls as well. Supplementary Table~\ref{stab:3} gives AUROC and sensitivity at re-estimated cut-offs. \pnl{e}, 40x immunofluorescence images of the six highest attribution cells from representative bladder, breast, colon, and prostate cancer patients with zero CK-gated cells in the sample. DAPI in blue, pan-cytokeratin in red, vimentin in white, and CD45/CD31 in green. Scale bars are 5~$\mu$m.}
\label{fig:4}
\end{figure}

\subsection*{Molecular profiling of high-attribution populations}

To resolve the lineage identity and genomic state of the circulating populations prioritized by SPECTRA, we cross-referenced mass-90 high-attribution cells against existing single-cell molecular datasets previously acquired from the same biobank slides, matching 48 cells (from 11 individuals) to multiplexed imaging mass cytometry\cite{gerdtsson2018,giesen2014} (IMC; 40-marker metal-conjugated antibody panel) and 108 cells (from 28 patients with prostate, breast, lung, or bladder cancer) to single-cell low-pass whole-genome sequencing\cite{malihi2020,baslan2012} (Fig.~\ref{fig:5}). Across the four positively emphasized populations represented in the IMC dataset (clusters 21, 17, 10, and 16), single-cell protein profiling distinguished epithelial from vascular and immune lineages (Fig.~\ref{fig:5}a). Profiled CK+, CD45/31$-$ cells in cluster 21 frequently expressed the epithelial adhesion molecule EpCAM (6 of 13 cells), whereas the most strongly up-weighted population, vimentin-high cluster 10 ($n=18$), rarely expressed EpCAM and instead included CD31+ endothelial-like cells and myeloid-marker-positive (CD14, CD68, or CD66b) cells (7 of 18 each), including two cells identified in non-cancer controls (Fig.~\ref{fig:5}a).

\begin{figure}[!tbp]
\centering
\fig{Figure5.pdf}
\caption{\textbf{Molecular characterization of high-attribution cells.} \pnl{a}, Single-cell proteomics via imaging mass cytometry (IMC) of 13 cells in cluster 21, 3 cells in cluster 17, 18 cells in cluster 10, and 14 cells in cluster 16. \pnl{b}, Single-cell low-pass whole-genome sequencing for copy-number variation analysis across cells in clusters 20, 21, 17, 10, and 16, with the epithelial clusters (20, 21, and 17) showing the highest proportion of aneuploid cells.}
\label{fig:5}
\end{figure}

Single-cell copy-number profiling across five high-attribution populations further separated cells with large-scale copy-number alterations from copy-number-neutral cells (Fig.~\ref{fig:5}b). Profiled cytokeratin-positive cells in clusters 20, 21, and 17 were predominantly aneuploid, harboring large-scale chromosomal alterations (28 of 31, 19 of 26, and 5 of 6 cells altered, respectively; Fig.~\ref{fig:5}b). By contrast, most profiled cells of vimentin-high cluster 10 were copy-number neutral (only 7 of 27 cells altered; 25.9\%), as were the CK+, CD45/31 dim, VIM+ cluster 16 (2 of 18 cells altered; 11.1\%), suggesting, in these cells, that aneuploidy tracks the combined phenotype rather than cytokeratin alone. Together, these profiled cells show that SPECTRA's high-attribution cells include both aneuploid, likely tumor-derived cells, and copy-number-neutral cells of endothelial and myeloid character.

\subsection*{Context-specific adaptation distinguishes cancer from high-risk clinical states}

A central challenge for blood-based cancer biology is identifying specific cancer signals in symptomatic patients, as opposed to genetically predisposed or environmentally influenced individuals who harbor benign inflammation or tissue remodeling rather than baseline physiology. To evaluate SPECTRA in these settings, we tested the locked pan-cancer ensemble on two reserved diagnostic cohorts excluded from primary model training: patients with lung cancer ($n=152$ training, $n=62$ test) versus individuals with a prior history of smoking undergoing diagnostic workup for lung nodules ($n=58$ training, $n=23$ test), and patients with breast cancer ($n=531$ training, $n=218$ test) versus cancer-free, high-risk BRCA-variant carriers ($n=14$ training, $n=8$ test).

Applied without modification, the pan-cancer SPECTRA showed reduced discrimination in both cohorts, with a test AUROC of 0.65 for lung cancer versus lung nodules and AUROC of 0.54 for breast cancer versus BRCA-variant carriers (Fig.~\ref{fig:6}a). Examining test-cohort score percentiles revealed that this drop in performance was attributable to high baseline scores in the diagnostic controls (Fig.~\ref{fig:6}b). Keeping the transformer frozen and refitting only the classification layer on the patient summary embedding improved diagnostic separation, raising test AUROC to 0.75 ($+0.10$, 95\% CI $-0.04$ to $+0.24$) for lung cancer versus lung nodules and to 0.89 ($+0.35$, 95\% CI $+0.10$ to $+0.62$, 8 test controls) for breast cancer versus BRCA carriers (Fig.~\ref{fig:6}a). These adapted models rest on few controls (14 training and 8 test BRCA-variant carriers) and are exploratory demonstrations that the pan-cancer representation can be re-read for task-specific purposes, not estimates of performance in high-risk populations. Other adaptation methods (retraining the last transformer block, or LoRA)\cite{hu2022} did not improve on decision layer adaptation (Supplementary Fig.~\ref{sfig:6}). Notably, this improvement was driven by suppressing elevated scores in high-risk controls while maintaining cancer patient scores (Fig.~\ref{fig:6}b).

\begin{figure}[!tbp]
\centering
\fig{Figure6.pdf}
\caption{\textbf{Context-specific adaptation distinguishes cancer from high-risk benign and genetic states.} \pnl{a}, Held-out test AUROC comparing the unchanged pan-cancer SPECTRA model (gray) against the adapted model with a refitted linear decision layer (teal) for lung cancer ($n=62$) versus lung nodule workup without confirmed cancer ($n=23$) and breast cancer ($n=218$) versus cancer-free BRCA-variant carriers ($n=8$). \pnl{b}, Mean SPECTRA score percentile within the held-out test set for cancers and diagnostic controls before and after adaptation. \pnl{c},\pnl{d}, Change in the share of positive cell-level attribution (percentage points) between the adapted and pan-cancer models across the most shifted clusters and all CK+ clusters combined. \pnl{e},\pnl{f}, 40x representative images of cells whose share of evidence rose (top) or dropped (bottom) after adaptation. DAPI in blue, pan-cytokeratin in red, vimentin in white, and CD45/CD31 in green. Scale bars are 10~$\mu$m.}
\label{fig:6}
\end{figure}

Comparing cell-level attribution before and after decision-layer refitting showed how SPECTRA re-weights circulating biology in distinct clinical contexts (Fig.~\ref{fig:6}c--f). In both clinical tasks, adaptation markedly decreased the reliance on large, vimentin-rich, CD45/CD31-variable cells (cluster 10), whose share of positive attribution fell by 3.7 percentage points in lung cancer and 10.1 percentage points in breast cancer (Fig.~\ref{fig:6}c,d). Of the 16 cells losing most weight, 14 belonged to cluster 10 (Fig.~\ref{fig:6}e,f, bottom rows), suggesting that these cells, which separate cancer from healthy donors, are less specific in comparisons with lung nodule work-up or BRCA carriers. Interestingly, in BRCA-variant carriers, the cluster 10 cells were also larger, with larger nuclei, and brighter for cytokeratin than in reference controls (patient-level AUROC 0.80--0.84, exploratory, 8 carriers), while we saw no such effect in lung-nodule workup patients. Conversely, each adapted classifier elevated context-specific circulating populations: lung cancer discrimination shifted attribution toward CD45/31 dim cells (cluster 5; $+3.7$ percentage points), DAPI-only marker-low cells (cluster 8; $+3.3$ percentage points), and a specific CK+, CD45/31$-$ population (cluster 21; $+2.8$ percentage points) while reducing overall CK+ share ($-2.6$ percentage points; Fig.~\ref{fig:6}c,e). Breast cancer discrimination against BRCA controls, on the other hand, increased weight on dual-positive CK+, CD45/31+ cells (clusters 18 and 19; $+1.6$ percentage points each) and total CK+ populations ($+3.1$ percentage points, Fig.~\ref{fig:6}d,f), indicating a shift towards dual-positive-like identification. Together, these findings indicate that the pan-cancer representation contains information that a small task-specific readout can use to separate these harder comparisons, and that doing so shifts weight away from vimentin-high cells.

% =====================================================================
\section*{Discussion}

SPECTRA, a transformer trained solely on patient-level diagnoses, reads enrichment-free images of rare circulating cells across 13 cancer types and traces each prediction back to the cells that drive it. Canonical aneuploid CK+, CD45/31$-$ cells were among the populations SPECTRA prioritized, yet counts of CK-gated candidate cells did not separate cancer patients from controls (AUROC 0.50), and removing them left discrimination unchanged. SPECTRA instead integrated contributions from multiple circulating populations: large, vimentin-high, copy-number-neutral cells expressing endothelial or myeloid markers carried the most attribution, alongside aneuploid CK+ cells and dual-positive CK+, CD45/31+ cells. SPECTRA also used differences in cell phenotype beyond the frequency of each population, suggesting that enumerating cell types or clusters alone does not capture the full biology. Workflows that enrich by surface markers or physical properties define in advance which cells are measured\cite{allard2004,nagrath2007,vona2000,ozkumur2013,sollier2014}, and would discard or overlook several of the populations SPECTRA weighted most heavily. Enrichment-free imaging thus captures cancer-associated information from a wider range of circulating cells than marker-based enumeration, encompassing both tumor-derived and likely host cells.

Among the profiled cells of the most heavily weighted population, the vimentin-high cells of cluster 10, most were copy-number-neutral and carried endothelial or myeloid markers, suggesting that much of the signal SPECTRA reads in blood arises from the host rather than the tumor. Circulating endothelial cells, tumor-associated macrophage-like cells and cells co-expressing epithelial and hematopoietic markers have each been linked to cancer\cite{bertolini2006,adams2014,ao2015,gast2018,higa2025,naghdloo2025caught}. Here, related populations were identified without prior hypotheses and tied directly to the model's decisions. A host-derived signal of this kind, if further confirmed, would explain why SPECTRA detected cancers in which no CK-gated cell was found, and why it performed comparably in early- and late-stage breast cancer. Defining these populations more precisely is a central next step. Because the archival slides were prepared from Streck-preserved blood, molecular profiling was confined to imaging mass cytometry and low-pass single-cell genome sequencing of a selected set of cells. Single-cell RNA sequencing of freshly collected samples will resolve the sub-populations within these clusters and their transcriptional states\cite{prabakar2024}.

The diagnostic control cohorts offer a further clue to the nature of this signal. The pan-cancer model scored most individuals undergoing lung-nodule work-up and most cancer-free BRCA-variant carriers as highly as patients with cancer. In BRCA-variant carriers, cluster 10 cells carried a larger share ($\sim$1.5x, 34\% vs 22--23\%) of the positive evidence than in either cancer patients or reference controls, and were larger and brighter for cytokeratin; in individuals undergoing lung-nodule work-up they did not carry a larger information share. Adaptation to these diagnostic controls reduced the weight placed on these cells. One interpretation is that in BRCA-variant carriers, vimentin-high cells report a vascular or myeloid response to tissue changes that precede cancer, such as field effects or premalignant changes in the breast tissue\cite{lim2009}, rather than to established tumors alone. These cells would then mark a biology shared by cancer and its precursors, explaining both their weight in the pan-cancer model and their reduced specificity in these comparisons. Longitudinal follow-up of such individuals, testing whether high scores or abundant cluster 10 cells precede a diagnosis, will be needed to distinguish this possibility from systemic inflammation or other benign processes.

These findings do not diminish the prognostic value of CK+ CTCs, whose enumeration predicts survival in metastatic breast, prostate and colorectal cancer, among others\cite{cristofanilli2004,debono2008,cohen2008}. Our study addresses a different question, namely how the circulating cellular compartment differs between patients with and without cancer, and in this setting CK+ counts carried no information. Whether the populations prioritized here also carry prognostic information can be tested within the same framework. Because SPECTRA learns from a patient-level objective, replacing the classification loss with a Cox proportional-hazards loss\cite{katzman2018} in outcome-annotated cohorts would allow attribution to identify the cells associated with survival in each indication.

SPECTRA is a proof of concept for reading cancer from the cellular compartment of blood, not a standalone detection or multi-cancer screening test. At 90\% specificity it detected 56.5\% of test cancer patients, and 39.7\% at 95\%. Population screening requires specificity of about 99\% or higher\cite{liu2020,klein2021}, which cannot be estimated from 44 controls. Part of the shortfall in sensitivity may be technical. SPECTRA analyzed approximately 1.5~mL of blood per patient, several-fold less than assays such as CellSearch, which examines 7.5~mL\cite{allard2004}. Under Poisson sampling\cite{tibbe2007}, a population present at one cell per 7.5~mL appears at least once in about 18\% of 1.5~mL samples but in about 63\% of 7.5~mL samples, so imaging more of the slides available from each draw offers a direct route to greater sensitivity. Beyond detection, SPECTRA makes the cellular compartment of blood measurable and interpretable at scale, which invites integration with cell-free DNA assays\cite{bettegowda2014,liu2020,cristiano2019,cohen2018}. Pairing the methylation, fragmentation and mutation signal of tumor DNA with the phenotypes of intact tumor and host cells would combine complementary views of the same disease. Early-fusion models that encode cell embeddings and cfDNA features\cite{bettegowda2014,liu2020,cristiano2019,cohen2018} jointly could use attribution to link specific cfDNA alterations to the circulating cell states that accompany them.

Several limitations qualify these results. Specificity was estimated from 44 held-out non-cancer controls, most of them healthy donors (90.9\% at the primary operating point; 95\% CI 78.8--96.4\%). Discrimination declined against clinically harder controls, namely individuals undergoing lung-nodule work-up and cancer-free BRCA-variant carriers, and recovered only partly after adaptation (Fig.~\ref{fig:6}). Controls were selected to have a median age similar to that of cancer patients, but most cancer and control samples were drawn from separate clinical studies. Our study-level analyses argue against site as the main source of signal, but excluding pre-analytical differences entirely will require controls collected alongside cases. Additionally, CTC enumeration in our study was conducted with thresholding and pre-trained cell-type classifiers, without expert review, so these counts approximate rather than reproduce validated CTC enumeration. All analyses used a single archival cohort processed with one assay, and we did not evaluate SPECTRA in an independent external validation cohort. Such validation, followed by prospective, multi-center studies with larger control populations, spanning healthy, benign, inflammatory and high-risk individuals, will be essential to define clinical utility and to test the biological interpretations proposed here.

Together, these results establish weakly supervised imaging of circulating cells as a scalable route to uncovering cancer biology in blood beyond canonical CK-positive CTCs, and lay a foundation for the next generation of cell-based liquid biopsies.

% =====================================================================
\section*{Methods}

\subsection*{Study design and participants}

This study used archival peripheral blood specimens from clinical studies and biobanks of patients with cancer and of individuals without cancer, processed with a common enrichment-free assay. All research was conducted in accordance with the principles of the Declaration of Helsinki. All participants, including the healthy donors provided by Epic Sciences, provided written informed consent prior to sample collection. All samples were de-identified to protect patient confidentiality. Institutional review boards approved all studies (see Supplementary Table~4 for IRB institutions and protocol numbers).

Peripheral blood samples were collected from 2,402 patients across 13 cancer types and 148 non-cancer controls: 106 healthy donors, 21 healthy women who donated blood following a cancer-negative mammogram, and 21 individuals with benign pancreatic disease sampled during pancreatic cancer workup (Table~\ref{tab:1}). Two further control populations were defined before model fitting as comparators for specific clinical questions and were excluded from general model fitting: 81 patients undergoing diagnostic workup for a lung nodule without a cancer diagnosis and 22 BRCA-variant carriers without a cancer diagnosis. Lung work-up controls were individuals with a history of smoking undergoing evaluation of lung nodules in whom cancer was not confirmed.

Clinical annotations were available for a set of breast cancer patients. Stage was recorded as early, late, or unknown. Early stage was taken from the recorded stage; late stage was defined by inclusion/exclusion criteria of the studies from which they were sampled (SWOG S1416\cite{rodler2023}: metastatic triple-negative; SWOG S1222\cite{moore2022}: metastatic HR+/HER2$-$). Of 218 test breast cancers, 50 were early, 98 late, and 70 unknown.

Patients were divided into a training partition (1,793 patients: 1,689 with cancer and 104 non-cancer controls) and a held-out test set partition (757 patients: 713 with cancer and 44 reference controls), allocating approximately 70\% of each cancer to training. The reserved diagnostic controls (lung and BRCA controls) were split in the same way for fine-tuning. Test patients were not used for model fitting, optimization, or selection, and were only scored once. Attribution, clustering, and the within-cluster analyses were performed on test patients after model locking and are descriptive.

\begin{table}[!tbp]
\centering
\caption{\textbf{Distribution of cancer types in the primary dataset.}}
\label{tab:1}
\begin{tabular}{lrrr}
\toprule
Cancer type & Train & Test & Total \\
\midrule
Prostate & 618 & 264 & 882 \\
Breast & 531 & 218 & 749 \\
Lung & 152 & 62 & 214 \\
Bladder & 121 & 52 & 173 \\
Upper tract urothelial carcinoma & 100 & 44 & 144 \\
Multiple myeloma & 48 & 23 & 71 \\
Pancreatic & 45 & 19 & 64 \\
Uveal melanoma & 29 & 12 & 41 \\
Bone tumors & 15 & 7 & 22 \\
Head and neck & 15 & 4 & 19 \\
Colon & 7 & 3 & 10 \\
Ovarian & 7 & 3 & 10 \\
Leptomeningeal disease of breast cancer & 1 & 2 & 3 \\
\midrule
\textbf{Total cancer patients} & \textbf{1,689} & \textbf{713} & \textbf{2,402} \\
\textbf{Non-cancer controls} & \textbf{104} & \textbf{44} & \textbf{148} \\
\bottomrule
\end{tabular}

\vspace{4pt}
\begin{minipage}{0.8\textwidth}\small
Reserved diagnostic controls (81 lung-nodule work-up, 22 BRCA-variant carriers) are not included in Table~\ref{tab:1}; see Fig.~\ref{fig:6}.
\end{minipage}
\end{table}

\subsection*{Sample processing, staining, and imaging}

Blood samples (mean blood volume of 8~mL) were obtained using 10~mL blood collection tubes (Cell-free DNA, Streck, La Vista, NE, USA) and were processed by the Convergent Science Institute for Cancer (CSI-Cancer) at the University of Southern California within 24--48 hours of collection, following a previously described procedure\cite{marrinucci2012}. In summary, ammonium chloride was used to lyse red blood cells and the remaining nucleated cell fraction was plated on specially designed glass slides (Marienfeld, Lauda, Baden-W\"urttemberg, Germany) at a concentration of approximately 2.5 million cells per slide. These slides were then stored at $-80\,^{\circ}$C for long-term preservation and rare cell analysis.

Up to two slides were analyzed per patient (two slides available for 2,467 of 2,550), totaling approximately five million analyzed cells per patient. Mean peripheral blood volume on each slide was 0.75~mL, with two slides representing a mean of 1.5~mL peripheral blood.

The slides were stained at room temperature using the IntelliPATH FLX\texttrademark\ autostainer (Biocare Medical LLC, Irvine, CA, USA), following a previously described method\cite{chai2021}. In summary, the samples were stained with a mouse IgG1 anti-human CD31:Alexa Fluor\textregistered\ 647 monoclonal antibody (clone: WM59, MCA1738A647, BioRad, Hercules, CA, USA) at a concentration of 2.5~$\mu$g/mL, and a goat anti-mouse IgG monoclonal Fab fragment (115-007-003, Jackson ImmunoResearch, West Grove, PA, USA) at a concentration of 100~$\mu$g/mL. The samples were then permeabilized using 100\% cold methanol and exposed to an antibody cocktail consisting of mouse IgG1/IgG2a anti-human cytokeratins (CKs) 1, 4, 5, 6, 8, 10, 13, 18, and 19 (clones: C-11, PCK-26, CY-90, KS-1A3, M20, A53-B/A2, C2562, Sigma, St.~Louis, MO, USA), mouse IgG1 anti-human CK 19 (clone: RCK108, GA61561-2, Dako, Carpinteria, CA, USA), mouse anti-human CD45:Alexa Fluor\textregistered\ 647 (clone: F10-89-4, MCA87A647, AbD Serotec, Raleigh, NC, USA), and rabbit IgG anti-human vimentin (Vim) (clone: D21H3, 9854BC, Cell Signaling, Danvers, MA, USA). Finally, the slides were incubated with Alexa Fluor\textregistered\ 555 goat anti-mouse IgG1 antibody (A21127, Invitrogen, Carlsbad, CA, USA) and 4$'$,6-diamidino-2-phenylindole (DAPI; D1306, ThermoFisher) and then mounted using a glycerol-based aqueous mounting medium. The samples were imaged using automated high-throughput fluorescence scanning microscopy with a 10x objective on custom in-house microscopes (pixel size: 0.59~$\mu$m $\times$ 0.59~$\mu$m).

Selected rare cells were re-imaged manually on a Zeiss Axioplan2 or Nikon Eclipse 80i with a 40x objective (pixel size: 0.18~$\mu$m $\times$ 0.18~$\mu$m). Re-imaged cells were used only for display and were not model inputs.

\subsection*{Single-cell segmentation, embedding, and rare-event detection}

After slide scanning, images were passed through a previously described embedding model. Cells were segmented utilizing a custom Cellpose cyto3 model\cite{naghdloo2025blue,stringer2021,stringer2025}. Each cell was represented by a 512-dimensional embedding produced by a frozen self-supervised encoder (SimCLR, ResNet-18 backbone)\cite{chen2020,he2016} from a $75\times75$ pixel, four-channel crop centered on the cell\cite{naghdloo2025blue}. The encoder was trained in a self-supervised manner on training-partition samples and never saw diagnosis labels. On each slide, an unsupervised outlier detector (ECOD) applied to the embeddings kept up to the 10,000 most anomalous cells as candidate rare events\cite{murgoitio2025,naghdloo2025blue,li2023}.

\subsection*{Cell filters}

Three filters, defined prior to model training, were applied to candidate rare events. First, a technical quality classifier, a frozen L2-regularized logistic regression on the cell embedding, gives an unusability score. The classifier was trained on 2,548 human-annotated usable/unusable rare events, and the Youden threshold\cite{youden1950} on the development set of 0.4651 was established for downstream use. Second, a leukocyte score was used to remove immune cells: cells with a white-blood-cell probability of greater than 0.9 were removed. Third, brightness: cells were kept only if at least one of the pan-CK, vimentin, or CD45/CD31 mean intensities was 5,000 or more (16-bit scale). A median of 1,440 cells were inputs to the SPECTRA model in the test partition (range 21 to 6,068).

\subsection*{SPECTRA model}

SPECTRA treats a patient as a set of unordered cells. Each cell embedding is standardized and projected linearly to width 256; a learned classification token (CLS) is prepended, and the sequence passes through six pre-norm transformer blocks\cite{xiong2020} (four attention heads, feed-forward expansion factor of 4 with GELU activation\cite{hendrycks2016}, dropout 0.4), with full bi-directional self-attention among all cells and the CLS token. There is no positional encoding. A final layer normalization of the CLS token and a linear classification layer give the patient logit, with probability defined after a sigmoid activation. The model has 4,870,913 trainable parameters. Embeddings were standardized with per-feature means and standard deviations estimated from the fitting patients only.

Models were trained from random initialization with binary cross-entropy loss against the cancer/control patient-level label. A batch size of 64 was used in training, with each minibatch containing 32 cancer patients and 32 controls. Cancer patients were sampled with cancer-type weights proportional to $n/\sqrt{\max(n,20)}$, where $n$ is the number of training patients of that cancer type, such that rare types were seen more often than their population-level prevalence. Optimization used the AdamW optimizer\cite{loshchilov2019} with a weight decay of 0.05, peak learning rate of $2\times10^{-4}$ with a 64-minibatch linear warmup and cosine decay\cite{loshchilov2017} to 5\% over a 30-epoch horizon, gradient clipping at 1.0, and 64 minibatches per epoch. Training ran for a fixed 15 epochs (960 updates) with no early stopping, with 15 epochs determined in cross-validation hyperparameter optimization.

\subsection*{Model selection}

Model hyperparameters (width, depth, number of attention heads, dropout, learning rate, and number of epochs) were chosen by three-fold cross-validation on the training partition, with folds balanced by cancer type. Two seeds were averaged per fold. The primary selection criterion was mean within-fold AUROC, with macro-averaged AUROC across nine cancer types with at least ten cases in the training split (eight cancer types with at least ten cases in the test split) as a secondary check. The final model is an ensemble of three six-layer models trained on all 1,793 training patients with seeds 17, 29, and 43, where the model score is the mean of the three probabilities. Ensemble agreement between seeds is shown in Supplementary Fig.~\ref{sfig:2}c; individual ensemble members reached test AUROCs of 0.74, 0.80, and 0.84.

We benchmarked additional cell pooling strategies derived from the multiple instance learning literature in three-fold cross-validation of the training set (Table~\ref{tab:2}). SPECTRA achieved higher mean fold AUROC than clustering-constrained attention-based multiple instance learning (CLAM\cite{lu2021}, single branch and multi-branch), attention-based multiple instance learning (ABMIL)\cite{ilse2018}, Additive MIL\cite{javed2022}, and mean pooling. All models used identical patient folds, input embeddings, cell filters, and normalization. Baselines used their published architectures with the same training budget and were not tuned separately. We additionally assessed whether the raw number of candidate events in a patient sample was predictive (count-based model), with this model showing predictive power near chance (AUROC 0.58--0.59). The count-based models applied a cubic-spline transformation to $\log(1+\text{events})$, followed by logistic regression with five coefficients and an intercept.

\begin{table}[!tbp]
\centering
\caption{\textbf{Matched three-fold cross-validation benchmark on the training partition} (identical folds, cell inputs, filters, normalization, and training budget).}
\label{tab:2}
\begin{tabular}{lrrr}
\toprule
Model & Parameters & Mean CV AUROC & Mean CV macro AUROC \\
\midrule
\textbf{SPECTRA} & \textbf{4,870,913} & \textbf{0.810} & \textbf{0.798} \\
CLAM-MB & 528,904 & 0.777 & 0.775 \\
CLAM-SB & 528,647 & 0.746 & 0.737 \\
ABMIL & 82,626 & 0.791 & 0.796 \\
Additive MIL & 82,626 & 0.757 & 0.748 \\
Mean pooling & 66,049 & 0.777 & 0.766 \\
Count-based model & 6 & 0.589 & 0.580 \\
\bottomrule
\end{tabular}
\end{table}

\subsection*{Test evaluation and endpoints}

The primary endpoint for model performance is the AUROC of the locked ensemble score for the 713 test cancer patients against the 44 test reference controls. Secondary endpoints are AUROCs per cancer type, macro-average AUROC over the eight cancer types with at least ten test patients, and sensitivity at fixed thresholds of 90\% and 95\% specificity (40 of 44 and 42 of 44 reference controls, respectively). Cancer types with fewer than ten test patients are reported descriptively. Endpoints using all of a patient's available slides (when more than two slides were available) gave similar results (Supplementary Fig.~\ref{sfig:2}e).

\subsection*{CK-gated and CTC-like cell counts}

To compare SPECTRA with conventional marker definitions in enrichment-based assays, we counted cells among each patient's model inputs, summed over the patient's analyzed slides. A cell was CK-gated if its mean CK intensity exceeded 40,000 and its CD45/CD31 intensity was below 10,000 (16-bit scale). CTC-like cells were those assigned by a frozen ten-cell-type classifier to its CTC, immune-like CTC, or platelet-coated CTC classes\cite{naghdloo2025blue}. These counts describe candidate cells, not confirmed tumor cells, and are intended to define two methods by which a CTC-enriched population can be defined. Scores across CK-gated strata were compared with the Kruskal--Wallis test. Correlations across clusters are Spearman. Within-cluster embedding AUROCs were tested against 2,000 label permutations with Benjamini--Hochberg correction. Composition versus SPECTRA used the paired stratified bootstrap above.

\subsection*{Batch effect estimation}

We assessed whether image acquisition contributed to SPECTRA performance. Primary AUROC was recomputed within subsets matched on scanner, scan year, scan date within 180 days, and DAPI-exposure quartile. Scores were also compared across scanners within cancers and within controls (Kruskal--Wallis). We further compared test scores against each control source (bootstrap AUROC, with each source left out); compared pancreatic cancer with pancreatic benign conditions from the same work-up study; compared scores among patients with the same cancer type across collection sites (Kruskal--Wallis test with $\varepsilon^2$ effect size, and pairwise AUROC); scored size-matched samples of leukocyte-classified candidate cells with the locked model; and fitted a logistic model of cancer status on SPECTRA score with DAPI exposure, scan year, and scanner as covariates. Each analysis was repeated on training split cross-validation out-of-fold scores. Collection sites were assigned per study from collection records (Supplementary Table~\ref{stab:2}). All of these analyses were performed after inspection of test results and are descriptive.

\subsection*{Statistical analysis}

AUROC 95\% confidence intervals were obtained by stratified patient bootstrapping\cite{efron1993}: cancers and controls were resampled separately, with 4,000 replicates. Macro AUROC was recomputed in each replicate. Paired differences, such as SPECTRA minus a marker count or fine-tuned minus pre-trained model, resampled the same patients for both scores. Sensitivity and specificity at the fixed thresholds have Wilson score intervals\cite{wilson1927}. Cluster abundance was compared between cancers and controls with Mann--Whitney U tests, with Benjamini--Hochberg correction\cite{benjamini1995} across clusters. Per-cancer AUROCs share the same 44 controls and are therefore correlated.

\subsection*{Cell attribution}

The contribution of each cell to a patient's score was estimated with AttnLRP, a layer-wise relevance propagation method for transformers. AttnLRP was chosen in a benchmark of attribution methods, which compared perturbation faithfulness on held-out training patients. Relevance was summed across embedding dimensions, giving one signed relevance per cell, and averaged over seeds in the ensemble. Two attributed sets were defined per patient: the 50 cells with the largest positive relevance (top-50), and the smallest set of cells carrying 90\% of the patient's total positive relevance (mass-90). Attribution was computed for all 757 test patients.

The mass-90 cells of all test patients (204,131 cells) were embedded by principal component analysis (PCA) to 50 components, fitted on these cells. UMAP\cite{mcinnes2018} was then used to display the PCA in two dimensions (30 nearest neighbors) and clustered with Leiden\cite{traag2019} (resolution 1), giving 23 clusters. Every model-input cell of each test patient was then assigned to the cluster of its 15 nearest mass-90 cells in the same 50-component space (98\% agreement with the atlas labels). Leiden runs with two other seeds gave 20--22 clusters (adjusted Rand index 0.56 and 0.51). Clusters were named by fixed rules from median marker intensities on the 16-bit scale: CK+ above 15,000; CK-dim 9,000--15,000; CD45/31 high above 35,000; CD45/31+ 15,000--35,000; CD45/31 dim 8,000--15,000; VIM+ above 10,000; DAPI bright above 20,000; DAPI dim below 6,000; large above 300 square pixels; and small below 140 square pixels.

For each patient, we computed each cluster's share of all model-input cells (abundance), and the emphasis of each cluster. Emphasis is $\log_2$ of a cluster's mean share among attributed cells over its mean share among all cells, in cancers. To ask whether cluster composition alone carries the signal, an L2 logistic regression ($C=0.1$, class-balanced) on centred log-ratio\cite{aitchison1982} cluster shares was fitted on the 1,793 training patients, whose cells were assigned to clusters as above, and scored once on the test patients. To ask whether cells within a cluster differ, each test patient with at least three cells in a cluster was represented by the mean 50-component embedding of those cells, and a logistic regression ($C=0.1$, class-balanced) was evaluated by $5\times5$ repeated stratified cross-validation within the test cohort, with 2,000 label permutations and Benjamini--Hochberg correction. The attribution atlas, the within-cluster analyses and the composition baseline use cells from the held-out test cohort and are exploratory.

For cluster 10, we computed 74 features with scikit-image (cell and nuclear shape; per-channel intensity statistics, radial distribution, nuclear-to-cytoplasmic ratio, and texture), averaged per patient. Groups were compared feature by feature (Mann--Whitney test with Benjamini--Hochberg correction) and jointly (L2-regularized logistic regression, $5\times5$ cross-validation, 200 label permutations).

\subsection*{Cell ablations}

We also tested how the SPECTRA score depends on the presence of specific sets of input cells. We removed the top-10, top-50, and mass-90 attributed cells, all CK-gated cells, all CTC-like cells, and all cells assigned to clusters 10, 20, and 21, in cancer patients and controls alike. We recorded AUROC with a stratified bootstrap 95\% CI and the paired difference from matched random removal; sensitivity and achieved specificity at the original cut-offs; and sensitivity at cut-offs re-estimated on the altered control scores to keep 40 of 44 (90\%) and 42 of 44 (95\%) controls below them. Removal changes the input distribution and the context of remaining cells, so these experiments show what the locked model depends on, not the causal role of removed cells.

\subsection*{Fine-tuning SPECTRA to specific clinical contexts}

Two comparisons used the reserved control populations: lung cancer against lung nodule workup (training 152 cancers, 58 controls; test 62 and 23) and breast cancer against BRCA-variant carriers (training 531 and 14; test 218 and 8). The locked ensemble was fine-tuned on each task's training patients. Three fine-tuning methods were used. First, the transformer was frozen and the final classifier was replaced by an L2 logistic regression ($C=0.1$, class-balanced) on the standardized CLS embedding of the patient. As a comparison, the final transformer block and classifier were retrained, with earlier blocks frozen, or low-rank adapters\cite{hu2022} (LoRA; query and value projections, rank 4, alpha 8) were trained together with the classifier. Both network adaptations used 5 epochs of 64 updates (AdamW, learning rate $5\times10^{-5}$, weight decay 0.05) without tuning (Supplementary Fig.~\ref{sfig:6}). CV estimates used the three-fold CV models as sources. Test AUROCs use the three-seed means, and gains over the pre-trained model have paired bootstrap intervals. Changes in cell attribution were measured as each cluster's share of positive relevance under the pre-trained and fine-tuned models. The decision-layer refit was chosen as the primary adaptation before inspection of test results in cross-validation, with LoRA and final-block refitting shown as illustrative comparisons.

\subsection*{Imaging mass cytometry and single-cell sequencing}

Existing imaging mass cytometry (IMC) and single-cell low-pass whole-genome sequencing were analyzed for further cluster interpretation. IMC was acquired on a Hyperion imaging system\cite{gerdtsson2018,giesen2014} (Standard BioTools) with panels of about 40 metal-conjugated antibodies. Each region was registered to its 10x or 40x DAPI frame by matching the IMC DNA signal (Ir191 and Ir193) to DAPI via an affine fit on matched nuclei. A region was accepted when at least half of its nuclei matched and the median residual was under three microns. Cells were segmented in OpenIMC (v1.1.3)\cite{tessone2025openimc} with the Cellpose nuclei model\cite{stringer2021} on the DNA channel. The masks were expanded to capture cell boundaries. Marker intensities were extracted with OpenIMC and arcsinh-scaled with a co-factor of 1. A marker was called positive if its total counts exceeded the background expectation over the same area (one-sided Poisson test, $p<0.001$), the mean was at least three times the background, and the mean was at least 0.3 counts per pixel. The IMC analysis is descriptive, covering 48 high-attribution cells from 11 individuals. Profiled cells were not explicitly selected for this study, but had been previously imaged by IMC or sequenced, and were matched based on their belonging to mass-90 attributed cells by slide position ($x$, $y$).

For sequencing, cells of interest were recovered from slides as previously described\cite{dago2014,malihi2020}. Each target cell was dislodged from the slide and aspirated into a micropipette on a micromanipulator (TransferMan NK2, Eppendorf). The cell was then deposited into a 0.2~mL PCR tube holding 1~$\mu$L Tris-EDTA buffer and kept at $-80\,^{\circ}$C until processing. Whole genome amplification used the GenomePlex single-cell kit (WGA4, Sigma-Aldrich) according to the manufacturer's instructions. Sequencing libraries were prepared from the purified amplification products with the NEBNext Ultra DNA Library Prep Kit for Illumina (New England Biolabs), including sonication and amplification. Library yield was measured by Qubit (Thermo Fisher). Fragment size was validated on an Agilent 2100 Bioanalyzer with the High Sensitivity DNA kit. Indexed libraries were pooled in batches of 50--70 at about 6~ng each, cleaned with AMPure XP beads (Beckman Coulter), and sequenced single-end at 50~bp on an Illumina NextSeq 500 or HiSeq 2500. Reads were trimmed by 30 bases at the $5'$ end to remove the WGA4 adapter, aligned to hg19 with Bowtie\cite{langmead2009}, sorted, and de-duplicated with SAMtools\cite{li2009}. Read counts were tallied across 5,000 genomic bins (0.6~Mb)\cite{baslan2012}. Each cell's profile was expressed as a $\log_2$ ratio to its median segment value. Noise was estimated per cell as 1.48 times the median absolute deviation of bin values from segment means. A segment was called altered when it spanned at least 17 consecutive bins ($\sim$10~Mb), its mean shifted by at least 0.2 from the median, and the shift exceeded six standard errors from noise. A cell was called altered when altered segments covered at least 5\% of the genome. In total, 108 mass-90 cells from 28 cancer patients were analyzed. These analyses are descriptive and describe cluster biology and the relative rate of aneuploidy within a cluster amongst sequenced cells.

\subsection*{Use of artificial intelligence}

Artificial intelligence tools (Anthropic Claude Opus 5.5 and OpenAI ChatGPT 6.1 Sol) were used for coding, data analysis, and text revision from a human-written source. The authors have verified all outputs.

% =====================================================================
\section*{Data availability}
Embeddings for outlier events and the cancer type of samples are made available in the BloodPAC Data Repository and in Zenodo.

\section*{Code availability}
All model architectures, analysis, and figure generation code is made available at \url{https://github.com/CSI-Cancer/spectra}.

\section*{Contributions}
Conceptualization: D.T., A.O., P.K.; Data curation: E.Q., S.S.; Formal analysis: D.T.; Investigation: D.T., Al.R., A.S., N.H., I.R., D.P., M.T.G., I.W., An.R., A.L., R.N., W.K., V.H.; Methodology: D.T., Al.R., A.S., J.M., A.O.; Project administration: J.H., P.K.; Resources: An.R., A.L., S.S.; Software: D.T., Al.R.; Validation: D.T., J.N., S.S., J.H., P.K.; Visualization: D.T.; Writing -- original draft: D.T.; Writing -- review and editing: all authors; Funding acquisition: J.M., A.O., J.H., P.K.

\section*{Acknowledgements}
We would like to acknowledge all of the patients who donated samples to our studies and biobank for research use, their families, and our clinical collaborators across institutions. We would also like to thank Dr.\ Amin Naghdloo for his thoughtful comments and conversations over the years.

\section*{Competing interests}
% to be completed

\section*{Funding}
% to be completed

\section*{Ethics declarations}
% to be completed

\section*{Reporting summary}
% to be completed

% =====================================================================
\begin{thebibliography}{99}
\small

\bibitem{bettegowda2014} Bettegowda, C. \textit{et al.} Detection of circulating tumor DNA in early- and late-stage human malignancies. \textit{Sci. Transl. Med.} \textbf{6}, 224ra24 (2014).
\bibitem{liu2020} Liu, M. C. \textit{et al.} Sensitive and specific multi-cancer detection and localization using methylation signatures in cell-free DNA. \textit{Ann. Oncol.} \textbf{31}, 745--759 (2020).
\bibitem{cristiano2019} Cristiano, S. \textit{et al.} Genome-wide cell-free DNA fragmentation in patients with cancer. \textit{Nature} \textbf{570}, 385--389 (2019).
\bibitem{dago2014} Dago, A. E. \textit{et al.} Rapid phenotypic and genomic change in response to therapeutic pressure in prostate cancer inferred by high content analysis of single circulating tumor cells. \textit{PLoS ONE} \textbf{9}, e101777 (2014).
\bibitem{gerdtsson2018} Gerdtsson, E. \textit{et al.} Multiplex protein detection on circulating tumor cells from liquid biopsies using imaging mass cytometry. \textit{Converg. Sci. Phys. Oncol.} \textbf{4}, 015002 (2018).
\bibitem{yu2013} Yu, M. \textit{et al.} Circulating breast tumor cells exhibit dynamic changes in epithelial and mesenchymal composition. \textit{Science} \textbf{339}, 580--584 (2013).
\bibitem{bertolini2006} Bertolini, F., Shaked, Y., Mancuso, P. \& Kerbel, R. S. The multifaceted circulating endothelial cell in cancer: towards marker and target identification. \textit{Nat. Rev. Cancer} \textbf{6}, 835--845 (2006).
\bibitem{adams2014} Adams, D. L. \textit{et al.} Circulating giant macrophages as a potential biomarker of solid tumors. \textit{Proc. Natl Acad. Sci. USA} \textbf{111}, 3514--3519 (2014).
\bibitem{ao2015} Ao, Z. \textit{et al.} Identification of cancer-associated fibroblasts in circulating blood from patients with metastatic breast cancer. \textit{Cancer Res.} \textbf{75}, 4681--4687 (2015).
\bibitem{gast2018} Gast, C. E. \textit{et al.} Cell fusion potentiates tumor heterogeneity and reveals circulating hybrid cells that correlate with stage and survival. \textit{Sci. Adv.} \textbf{4}, eaat7828 (2018).
\bibitem{higa2025} Higa, N. \textit{et al.} Simultaneous expression of epithelial and immune cell markers in circulating tumor cells identified in patients with stage 4 breast cancer. \textit{Commun. Med.} \textbf{5}, 309 (2025).
\bibitem{naghdloo2025caught} Naghdloo, A. \textit{et al.} Caught in the act: tumor-immune interactions in circulation of patients with immune marker positive circulating tumor cells. \textit{Cancers} \textbf{17}, 3667 (2025).
\bibitem{allard2004} Allard, W. J. \textit{et al.} Tumor cells circulate in the peripheral blood of all major carcinomas but not in healthy subjects or patients with nonmalignant diseases. \textit{Clin. Cancer Res.} \textbf{10}, 6897--6904 (2004).
\bibitem{nagrath2007} Nagrath, S. \textit{et al.} Isolation of rare circulating tumour cells in cancer patients by microchip technology. \textit{Nature} \textbf{450}, 1235--1239 (2007).
\bibitem{vona2000} Vona, G. \textit{et al.} Isolation by size of epithelial tumor cells: a new method for the immunomorphological and molecular characterization of circulating tumor cells. \textit{Am. J. Pathol.} \textbf{156}, 57--63 (2000).
\bibitem{ozkumur2013} Ozkumur, E. \textit{et al.} Inertial focusing for tumor antigen-dependent and -independent sorting of rare circulating tumor cells. \textit{Sci. Transl. Med.} \textbf{5}, 179ra47 (2013).
\bibitem{sollier2014} Sollier, E. \textit{et al.} Size-selective collection of circulating tumor cells using Vortex technology. \textit{Lab Chip} \textbf{14}, 63--77 (2014).
\bibitem{marrinucci2012} Marrinucci, D. \textit{et al.} Fluid biopsy in patients with metastatic prostate, pancreatic and breast cancers. \textit{Phys. Biol.} \textbf{9}, 016003 (2012).
\bibitem{murgoitio2025} Murgoitio-Esandi, J. \textit{et al.} Unsupervised detection of rare events in liquid biopsy assays. \textit{npj Precis. Oncol.} \textbf{9}, 225 (2025).
\bibitem{naghdloo2025blue} Naghdloo, A. \textit{et al.} Representation learning enables robust single cell phenotyping in whole slide liquid biopsy imaging. \textit{Sci. Rep.} \textbf{15}, 36589 (2025).
\bibitem{bai2024} Bai, L. \textit{et al.} Longitudinal tracking of circulating rare events in the liquid biopsy of stage III--IV non-small cell lung cancer patients. \textit{Discov. Oncol.} \textbf{15}, 142 (2024).
\bibitem{resnick2024} Resnick, K. \textit{et al.} Circulation of rare events in the liquid biopsy for early detection of lung mass lesions. \textit{Thorac. Cancer} \textbf{15}, 2100--2109 (2024).
\bibitem{malihi2020} Malihi, P. D. \textit{et al.} Single-cell circulating tumor cell analysis reveals genomic instability as a distinctive feature of aggressive prostate cancer. \textit{Clin. Cancer Res.} \textbf{26}, 4143--4153 (2020).
\bibitem{shishido2025} Shishido, S. N. \textit{et al.} Characterizing circulating rare cells in peripheral blood for detecting and monitoring multiple myeloma and precursor states. \textit{npj Precis. Oncol.} \textbf{9}, 388 (2025).
\bibitem{campanella2019} Campanella, G. \textit{et al.} Clinical-grade computational pathology using weakly supervised deep learning on whole slide images. \textit{Nat. Med.} \textbf{25}, 1301--1309 (2019).
\bibitem{lu2021} Lu, M. Y. \textit{et al.} Data-efficient and weakly supervised computational pathology on whole-slide images. \textit{Nat. Biomed. Eng.} \textbf{5}, 555--570 (2021).
\bibitem{courtiol2019} Courtiol, P. \textit{et al.} Deep learning-based classification of mesothelioma improves prediction of patient outcome. \textit{Nat. Med.} \textbf{25}, 1519--1525 (2019).
\bibitem{skrede2020} Skrede, O.-J. \textit{et al.} Deep learning for prediction of colorectal cancer outcome: a discovery and validation study. \textit{Lancet} \textbf{395}, 350--360 (2020).
\bibitem{vaswani2017} Vaswani, A. \textit{et al.} Attention is all you need. \textit{Adv. Neural Inf. Process. Syst.} \textbf{30}, 5998--6008 (2017).
\bibitem{lee2019} Lee, J. \textit{et al.} Set Transformer: a framework for attention-based permutation-invariant neural networks. \textit{Proc. 36th Int. Conf. Mach. Learn.} \textbf{97}, 3744--3753 (2019).
\bibitem{shao2021} Shao, Z. \textit{et al.} TransMIL: Transformer based correlated multiple instance learning for whole slide image classification. \textit{Adv. Neural Inf. Process. Syst.} \textbf{34}, 2136--2147 (2021).
\bibitem{wagner2023} Wagner, S. J. \textit{et al.} Transformer-based biomarker prediction from colorectal cancer histology: a large-scale multicentric study. \textit{Cancer Cell} \textbf{41}, 1650--1661.e4 (2023).
\bibitem{bach2015} Bach, S. \textit{et al.} On pixel-wise explanations for non-linear classifier decisions by layer-wise relevance propagation. \textit{PLoS ONE} \textbf{10}, e0130140 (2015).
\bibitem{achtibat2024} Achtibat, R. \textit{et al.} AttnLRP: attention-aware layer-wise relevance propagation for transformers. \textit{Proc. 41st Int. Conf. Mach. Learn.} \textbf{235}, 135--168 (2024).
\bibitem{traag2019} Traag, V. A., Waltman, L. \& van Eck, N. J. From Louvain to Leiden: guaranteeing well-connected communities. \textit{Sci. Rep.} \textbf{9}, 5233 (2019).
\bibitem{mcinnes2018} McInnes, L., Healy, J. \& Melville, J. UMAP: uniform manifold approximation and projection for dimension reduction. Preprint at \url{https://arxiv.org/abs/1802.03426} (2018).
\bibitem{giesen2014} Giesen, C. \textit{et al.} Highly multiplexed imaging of tumor tissues with subcellular resolution by mass cytometry. \textit{Nat. Methods} \textbf{11}, 417--422 (2014).
\bibitem{baslan2012} Baslan, T. \textit{et al.} Genome-wide copy number analysis of single cells. \textit{Nat. Protoc.} \textbf{7}, 1024--1041 (2012).
\bibitem{hu2022} Hu, E. J. \textit{et al.} LoRA: low-rank adaptation of large language models. In \textit{Proc. Int. Conf. Learn. Represent.} (2022).
\bibitem{prabakar2024} Prabakar, R. K., Schmidt, M. J., Kuhn, P. \& Hicks, J. GeTMoR: simultaneous genomic, transcriptomic, and morphological profiling of rare single cells. Preprint at bioRxiv \url{https://doi.org/10.1101/2024.09.29.615279} (2024).
\bibitem{slaughter1953} Slaughter, D. P., Southwick, H. W. \& Smejkal, W. Field cancerization in oral stratified squamous epithelium; clinical implications of multicentric origin. \textit{Cancer} \textbf{6}, 963--968 (1953).
\bibitem{steiling2008} Steiling, K., Ryan, J., Brody, J. S. \& Spira, A. The field of tissue injury in the lung and airway. \textit{Cancer Prev. Res.} \textbf{1}, 396--403 (2008).
\bibitem{lim2009} Lim, E. \textit{et al.} Aberrant luminal progenitors as the candidate target population for basal tumor development in BRCA1 mutation carriers. \textit{Nat. Med.} \textbf{15}, 907--913 (2009).
\bibitem{cristofanilli2004} Cristofanilli, M. \textit{et al.} Circulating tumor cells, disease progression, and survival in metastatic breast cancer. \textit{N. Engl. J. Med.} \textbf{351}, 781--791 (2004).
\bibitem{debono2008} de Bono, J. S. \textit{et al.} Circulating tumor cells predict survival benefit from treatment in metastatic castration-resistant prostate cancer. \textit{Clin. Cancer Res.} \textbf{14}, 6302--6309 (2008).
\bibitem{cohen2008} Cohen, S. J. \textit{et al.} Relationship of circulating tumor cells to tumor response, progression-free survival, and overall survival in patients with metastatic colorectal cancer. \textit{J. Clin. Oncol.} \textbf{26}, 3213--3221 (2008).
\bibitem{katzman2018} Katzman, J. L. \textit{et al.} DeepSurv: personalized treatment recommender system using a Cox proportional hazards deep neural network. \textit{BMC Med. Res. Methodol.} \textbf{18}, 24 (2018).
\bibitem{klein2021} Klein, E. A. \textit{et al.} Clinical validation of a targeted methylation-based multi-cancer early detection test using an independent validation set. \textit{Ann. Oncol.} \textbf{32}, 1167--1177 (2021).
\bibitem{tibbe2007} Tibbe, A. G. J., Miller, M. C. \& Terstappen, L. W. M. M. Statistical considerations for enumeration of circulating tumor cells. \textit{Cytometry A} \textbf{71}, 154--162 (2007).
\bibitem{cohen2018} Cohen, J. D. \textit{et al.} Detection and localization of surgically resectable cancers with a multi-analyte blood test. \textit{Science} \textbf{359}, 926--930 (2018).
\bibitem{rodler2023} Rodler, E. \textit{et al.} Cisplatin with veliparib or placebo in metastatic triple-negative breast cancer and BRCA mutation-associated breast cancer (S1416): a randomised, double-blind, placebo-controlled, phase 2 trial. \textit{Lancet Oncol.} \textbf{24}, 162--174 (2023).
\bibitem{moore2022} Moore, H. C. F. \textit{et al.} A randomized trial of fulvestrant, everolimus, and anastrozole for the front-line treatment of patients with advanced hormone receptor-positive breast cancer, SWOG S1222. \textit{Clin. Cancer Res.} \textbf{28}, 611--617 (2022).
\bibitem{chai2021} Chai, S. \textit{et al.} Platelet-coated circulating tumor cells are a predictive biomarker in patients with metastatic castrate-resistant prostate cancer. \textit{Mol. Cancer Res.} \textbf{19}, 2036--2045 (2021).
\bibitem{stringer2021} Stringer, C., Wang, T., Michaelos, M. \& Pachitariu, M. Cellpose: a generalist algorithm for cellular segmentation. \textit{Nat. Methods} \textbf{18}, 100--106 (2021).
\bibitem{stringer2025} Stringer, C. \& Pachitariu, M. Cellpose3: one-click image restoration for improved cellular segmentation. \textit{Nat. Methods} \textbf{22}, 592--599 (2025).
\bibitem{chen2020} Chen, T., Kornblith, S., Norouzi, M. \& Hinton, G. A simple framework for contrastive learning of visual representations. \textit{Proc. 37th Int. Conf. Mach. Learn.} \textbf{119}, 1597--1607 (2020).
\bibitem{he2016} He, K., Zhang, X., Ren, S. \& Sun, J. Deep residual learning for image recognition. In \textit{Proc. IEEE Conf. Comput. Vis. Pattern Recognit.} 770--778 (2016).
\bibitem{li2023} Li, Z. \textit{et al.} ECOD: unsupervised outlier detection using empirical cumulative distribution functions. \textit{IEEE Trans. Knowl. Data Eng.} \textbf{35}, 12181--12193 (2023).
\bibitem{youden1950} Youden, W. J. Index for rating diagnostic tests. \textit{Cancer} \textbf{3}, 32--35 (1950).
\bibitem{xiong2020} Xiong, R. \textit{et al.} On layer normalization in the transformer architecture. \textit{Proc. 37th Int. Conf. Mach. Learn.} \textbf{119}, 10524--10533 (2020).
\bibitem{hendrycks2016} Hendrycks, D. \& Gimpel, K. Gaussian error linear units (GELUs). Preprint at \url{https://arxiv.org/abs/1606.08415} (2016).
\bibitem{loshchilov2019} Loshchilov, I. \& Hutter, F. Decoupled weight decay regularization. In \textit{Proc. Int. Conf. Learn. Represent.} (2019).
\bibitem{loshchilov2017} Loshchilov, I. \& Hutter, F. SGDR: stochastic gradient descent with warm restarts. In \textit{Proc. Int. Conf. Learn. Represent.} (2017).
\bibitem{ilse2018} Ilse, M., Tomczak, J. M. \& Welling, M. Attention-based deep multiple instance learning. \textit{Proc. 35th Int. Conf. Mach. Learn.} \textbf{80}, 2127--2136 (2018).
\bibitem{javed2022} Javed, S. A. \textit{et al.} Additive MIL: intrinsically interpretable multiple instance learning for pathology. \textit{Adv. Neural Inf. Process. Syst.} \textbf{35} (2022).
\bibitem{efron1993} Efron, B. \& Tibshirani, R. J. \textit{An Introduction to the Bootstrap} (Chapman \& Hall, 1993).
\bibitem{wilson1927} Wilson, E. B. Probable inference, the law of succession, and statistical inference. \textit{J. Am. Stat. Assoc.} \textbf{22}, 209--212 (1927).
\bibitem{benjamini1995} Benjamini, Y. \& Hochberg, Y. Controlling the false discovery rate: a practical and powerful approach to multiple testing. \textit{J. R. Stat. Soc. B} \textbf{57}, 289--300 (1995).
\bibitem{adebayo2018} Adebayo, J. \textit{et al.} Sanity checks for saliency maps. \textit{Adv. Neural Inf. Process. Syst.} \textbf{31}, 9505--9515 (2018).
\bibitem{aitchison1982} Aitchison, J. The statistical analysis of compositional data. \textit{J. R. Stat. Soc. B} \textbf{44}, 139--177 (1982).
\bibitem{tessone2025openimc} Tessone, D. OpenIMC (v1.1.3). GitHub \url{https://github.com/dean-tessone/OpenIMC} (2025).
\bibitem{langmead2009} Langmead, B., Trapnell, C., Pop, M. \& Salzberg, S. L. Ultrafast and memory-efficient alignment of short DNA sequences to the human genome. \textit{Genome Biol.} \textbf{10}, R25 (2009).
\bibitem{li2009} Li, H. \textit{et al.} The Sequence Alignment/Map format and SAMtools. \textit{Bioinformatics} \textbf{25}, 2078--2079 (2009).
\end{thebibliography}

% =====================================================================
\clearpage
\nolinenumbers
\section*{Supplementary Information}

\setcounter{figure}{0}
\setcounter{table}{0}
\renewcommand{\figurename}{Supplementary Fig.}
\renewcommand{\tablename}{Supplementary Table}

\begin{figure}[!htb]
\centering
\fig{SupplementaryFig1.pdf}
\caption{\textbf{Cell inputs to SPECTRA.} \pnl{a}, Cells are embedded with a pre-trained encoder, BLUE (Biomarker-driven liquid biopsy unified encoder) v2, which was pre-trained on 3.8 million rare events from 507 patients across cancer types. \pnl{b}, Cell cascades into SPECTRA. After rare event detection, 20,000 cells are selected per patient ($\sim$0.5\% of cells). An image QC filter brings that down to a median of 17,126 cells/patient, and subsequent removal of immune cells and low-intensity outliers resolves to a median of 1,426 cells/patient. \pnl{c}, SPECTRA sees cells of various classes based on pre-trained classifiers. Logistic regression of cell-type predicts a median of 605 CTC-like cells per patient, 478 cells that fit a large-nuclei phenotype with no marker expression, 24 endothelial-like cells, 14 megakaryocyte-like cells, 1 fibroblast-like cell, and 2 CK-gated cells.}
\label{sfig:1}
\end{figure}

\begin{figure}[p]
\centering
\fig{SupplementaryFig2.pdf}
\caption{\textbf{SPECTRA technical ablations.} \pnl{a}, Comparison of SPECTRA (6-layer, CLS-supervised Transformer) to other multiple-instance learning (MIL) architectures in cross-validation. ABMIL, attention-based multiple instance learning; CLAM-MB and CLAM-SB, clustering-constrained attention-based multiple-instance learning (multi-branch and single-branch, respectively). \pnl{b}, Difference in performance to SPECTRA for each MIL architecture in each fold. \pnl{c}, Ensemble performance assessing how many cancers and controls are correctly classified in each of three seeds. \pnl{d}, Spearman correlation between the number of input cells to SPECTRA and model score. \pnl{e}, Assessing the effect of the exact set of 2 slides used per patient; the two slides were selected at random (primary group). For patients with more than 2 available slides ($n=8$), we assess selecting a different set of slides (same draw and tube) or providing all slides (all eligible groups). \pnl{f}, Technical variate effects on SPECTRA performance (comparing cancer and controls from the same scanner instrument, same year, $\le$180 days apart, and similar exposure settings).}
\label{sfig:2}
\end{figure}

\begin{figure}[p]
\centering
\fig{SupplementaryFig3.pdf}
\caption{\textbf{Study, collection-site and acquisition checks.} \pnl{a}, Test scores of cancer patients (violin, $n=713$) and of reference controls by source (points; $n$ below); numbers above, AUROC of all cancers against each source; bar, median; dashed line, locked cut-off at 90\% specificity. \pnl{b}, AUROC (95\% CI) for all cancers against all reference controls, and for pancreatic cancer against benign pancreatic disease within one work-up study and at one site; filled, test set; open, training cross-validation (out-of-fold); numbers, test cancers/controls. \pnl{c}, Test scores of patients with hormone-receptor-positive breast cancer and with lung cancer by collection site; numbers, patients; bar, median; $P$ and $\varepsilon^2$, Kruskal--Wallis test across sites and its effect size. No site-associated difference was detected, which does not establish independence from site; site is also correlated with disease stage and cohort purpose. \pnl{d}, AUROC (95\% CI) for distinguishing patients with the same cancer type collected at two sites (0.5, indistinguishable); filled, test; open, training out-of-fold. \pnl{e}, Test AUROC (95\% CI) of the locked model applied to its real input cells and to size-matched bags of candidate cells classified as leukocytes, which the model never receives. Descriptive analyses performed after test inspection.}
\label{sfig:sites}
\end{figure}

\begin{figure}[p]
\centering
\fig{SupplementaryFig4.pdf}
\caption{\textbf{Assessing the learning of relevant cells in SPECTRA.} \pnl{a}, Random re-initialization of each layer in the model, step-wise cumulatively, comparing the attribution of cells layer-wise\cite{adebayo2018}. We see that significant performance is learned in layers 2 and 3 of the transformer, and the most information is pooled at the classifier level. \pnl{b}, Seed-wise agreement of emphasis compared to the ensemble. \pnl{c}, Positive predictive mass in the test set, assessing the number of cells needed to capture any percentage of predictive evidence.}
\label{sfig:3}
\end{figure}

\begin{figure}[p]
\centering
\fig{SupplementaryFig5.pdf}
\caption{\textbf{Gallery of 40x microscopy of cells in each of 23 recurrent populations identified by SPECTRA.} DAPI in blue, pan-cytokeratin in red, vimentin in white, and CD45/CD31 in green. Scale bars are 10~$\mu$m.}
\label{sfig:4}
\end{figure}

\begin{figure}[p]
\centering
\fig{SupplementaryFig6.pdf}
\caption{\textbf{Continuous phenotype versus clustering.} \pnl{a}, Within each cluster, comparison of the AUROC (cross-validation in the test set) for predicting cancer versus control from the abundance of the cluster and from the mean embedding of a patient's cells within the cluster. \pnl{b}, Scatter plot comparison of abundance and embedding performance (cross-validation in the test set).}
\label{sfig:5}
\end{figure}

\begin{figure}[p]
\centering
\fig{SupplementaryFig7.pdf}
\caption{\textbf{Comparing different fine-tuning methods of SPECTRA.} \pnl{a}, Test AUROC performance of the original pan-cancer SPECTRA, final-layer refitting, last transformer block and decision-layer retraining, and LoRA and decision-layer retraining in lung cancer vs benign lung nodules (left) and breast cancer vs BRCA carriers (right). \pnl{b}, Training CV AUROCs for the same comparisons to assess train-test gaps and overfitting.}
\label{sfig:6}
\end{figure}

\clearpage
\begin{table}[!tbp]
\centering
\caption{\textbf{SPECTRA performance by cancer type on the held-out test set.} Each cancer type versus the same 44 non-cancer controls; AUROCs for different cancer types share these controls and are therefore correlated. AUROC 95\% CIs from a stratified patient bootstrap (4,000 replicates); not given for types with fewer than 10 patients, which are descriptive. Sensitivity at the common cut-offs set on the pooled ROC at 90\% and 95\% specificity (achieved 40/44 and 42/44), with Wilson 95\% intervals.}
\label{stab:1}
\small
\begin{tabular}{lrlll}
\toprule
Cancer type & $n$ & AUROC (95\% CI) & Sens.\ at 90\% spec.\ (95\% CI) & Sens.\ at 95\% spec.\ (95\% CI) \\
\midrule
All cancers & 713 & 0.82 (0.75--0.89) & 56.5\% (403/713; 52.9--60.1) & 39.7\% (283/713; 36.2--43.3) \\
Prostate & 264 & 0.81 (0.73--0.87) & 51.5\% (136/264; 45.5--57.5) & 36.0\% (95/264; 30.4--41.9) \\
Breast & 218 & 0.84 (0.76--0.90) & 58.7\% (128/218; 52.1--65.0) & 39.4\% (86/218; 33.2--46.1) \\
Lung & 62 & 0.90 (0.83--0.96) & 82.3\% (51/62; 71.0--89.8) & 54.8\% (34/62; 42.5--66.6) \\
Bladder & 52 & 0.82 (0.73--0.90) & 53.8\% (28/52; 40.5--66.7) & 44.2\% (23/52; 31.6--57.7) \\
UTUC & 44 & 0.83 (0.73--0.90) & 54.5\% (24/44; 40.1--68.3) & 45.5\% (20/44; 31.7--59.9) \\
Myeloma \& related & 23 & 0.76 (0.64--0.88) & 47.8\% (11/23; 29.2--67.0) & 30.4\% (7/23; 15.6--50.9) \\
Pancreatic & 19 & 0.88 (0.78--0.96) & 68.4\% (13/19; 46.0--84.6) & 47.4\% (9/19; 27.3--68.3) \\
Uveal melanoma & 12 & 0.72 (0.55--0.87) & 25.0\% (3/12; 8.9--53.2) & 25.0\% (3/12; 8.9--53.2) \\
Bone tumors & 7 & 0.81 & 71.4\% (5/7; 35.9--91.8) & 42.9\% (3/7; 15.8--75.0) \\
Head and neck & 4 & 0.65 & 25.0\% (1/4; 4.6--69.9) & 0.0\% (0/4; 0.0--49.0) \\
Ovarian & 3 & 0.52 & 0.0\% (0/3; 0.0--56.1) & 0.0\% (0/3; 0.0--56.1) \\
Colon & 3 & 0.91 & 66.7\% (2/3; 20.8--93.9) & 66.7\% (2/3; 20.8--93.9) \\
Leptomeningeal & 2 & 0.77 & 50.0\% (1/2; 9.5--90.5) & 50.0\% (1/2; 9.5--90.5) \\
\bottomrule
\end{tabular}
\end{table}

\clearpage
\begin{table}[p]
\centering
\caption{\textbf{Studies and collection sites.} Patients by cancer type and control source in the training and held-out test partitions, with the number of contributing studies and their collection sites (from laboratory records; multi-site trials are listed as trial sites; sites ordered by number of patients). Reference controls (Scripps normal donors, Epic Sciences, Cedars-Sinai INTC17-C and benign pancreatic work-up) form the primary comparison; reserved controls were used only for the adaptation analyses (Fig.~\ref{fig:6}). JN-family lung patients were collected through USC and Billings Clinic without a per-patient record and are listed under USC. Per-study counts, scan years, scanners and tube types are in Supplementary Data 1. NR, not recorded.}
\label{stab:2}
\footnotesize
\setlength{\tabcolsep}{4pt}
\begin{tabular}{L{3.5cm} C{1.1cm} L{8.2cm} C{1.9cm}}
\toprule
Group & Studies & Collection sites & Patients (train / test) \\
\midrule
Prostate & 8 & MD Anderson, Morehouse/Grady, UCLA, City of Hope, Johns Hopkins, UPenn, UAB & 618 / 264 \\
Breast & 13 & Billings Clinic, University of Colorado, Duke, Walter Reed, Anne Arundel, Cedars-Sinai, City of Hope, USC, MD Anderson; SWOG S1416 and NRG-BR002 trial sites; one study NR & 531 / 218 \\
Lung & 10 & Stanford, USC, MSKCC, Manchester & 152 / 62 \\
Bladder & 4 & USC, Johns Hopkins, UCSD & 121 / 52 \\
UTUC & 1 & USC & 100 / 44 \\
Myeloma and related states & 1 & MD Anderson & 48 / 23 \\
Pancreatic & 3 & Cedars-Sinai, Mayo Clinic & 45 / 19 \\
Uveal melanoma & 1 & USC & 29 / 12 \\
Bone tumors & 1 & CHLA & 15 / 7 \\
Head and neck & 1 & USC & 15 / 4 \\
Colon & 1 & Gilead trial sites & 7 / 3 \\
Ovarian & 1 & City of Hope & 7 / 3 \\
Leptomeningeal disease & 1 & City of Hope & 1 / 2 \\
Reference controls & 4 & Scripps, Epic Sciences, Cedars-Sinai & 104 / 44 \\
Lung-nodule work-up (reserved) & 3 & Stanford, USC & 58 / 23 \\
BRCA-variant carriers (reserved) & 1 & UChicago & 14 / 8 \\
\bottomrule
\end{tabular}
\end{table}
\clearpage

\begin{table}[p]
\centering
\caption{\textbf{In silico cell removal on the locked model, and CK-gate thresholds.} \textbf{a}, Each cell set was removed from every held-out test patient (713 cancers, 44 controls; cancers and controls alike) and the locked three-model ensemble re-scored; matched random removal deletes the same number of the patient's other cells (mean of five draws). AUROC 95\% CIs and differences from stratified patient bootstrap (4,000 replicates). Sensitivity is given at the original cut-off (with the number of controls below it) and at a cut-off re-estimated on the altered control scores to keep 40 of 44 controls below. Removal alters the input distribution, so these results show what the model output depends on, not the causal role of the removed cells. \textbf{b}, AUROC of per-patient CK-gated cell counts across nine predefined gate thresholds (16-bit intensities).}
\label{stab:3}
\footnotesize
\setlength{\tabcolsep}{4pt}
\textbf{a}\par\smallskip
\begin{tabular}{L{3.1cm} C{1.4cm} L{2.3cm} C{1.5cm} L{3.4cm} L{2.4cm} C{1.7cm}}
\toprule
Cells removed & Median per cancer patient & AUROC (95\% CI) & AUROC, matched random & Difference vs random (95\% CI) & Sensitivity at original 90\% cut-off (controls below) & Sensitivity at re-estimated 90\% cut-off \\
\midrule
None (intact model) & -- & 0.82 (0.75--0.89) & -- & -- & 56.5\% (40/44) & 56.5\% \\
CK-gated cells & 2 & 0.82 (0.75--0.88) & 0.83 & -0.002 (-0.009 to +0.004) & 56.2\% (40/44) & 56.4\% \\
All CTC-like cells & 618 & 0.73 (0.65--0.80) & 0.80 & -0.067 (-0.157 to +0.025) & 59.3\% (36/44) & 32.1\% \\
Top-10 attributed cells & 10 & 0.84 (0.78--0.89) & 0.82 & +0.013 (-0.023 to +0.053) & 28.9\% (43/44) & 51.6\% \\
Top-50 attributed cells & 50 & 0.82 (0.76--0.88) & 0.83 & -0.004 (-0.065 to +0.069) & 14.9\% (44/44) & 60.4\% \\
Mass-90 attributed cells & 220 & 0.78 (0.71--0.84) & 0.82 & -0.039 (-0.105 to +0.028) & 0.6\% (44/44) & 45.7\% \\
All cluster 10 cells & 32 & 0.83 (0.77--0.89) & 0.83 & +0.004 (-0.035 to +0.046) & 42.1\% (42/44) & 60.4\% \\
All cluster 20 cells & 23 & 0.82 (0.74--0.88) & 0.82 & -0.006 (-0.018 to +0.004) & 54.1\% (39/44) & 52.5\% \\
All cluster 21 cells & 31 & 0.82 (0.75--0.89) & 0.83 & -0.005 (-0.025 to +0.015) & 56.1\% (39/44) & 48.9\% \\
\bottomrule
\end{tabular}

\bigskip
\textbf{b}\par\smallskip
\begin{tabular}{lll}
\toprule
CK gate (CK $>$ / CD45/CD31 $<$) & AUROC (95\% CI) & Controls with $\geq$1 gated cell (of 44) \\
\midrule
30k / 5k & 0.52 (0.44--0.61) & 30 \\
30k / 10k & 0.54 (0.45--0.62) & 39 \\
30k / 15k & 0.56 (0.48--0.64) & 43 \\
40k / 5k & 0.51 (0.44--0.59) & 25 \\
40k / 10k (primary) & 0.50 (0.42--0.59) & 32 \\
40k / 15k & 0.53 (0.45--0.60) & 37 \\
50k / 5k & 0.50 (0.43--0.56) & 14 \\
50k / 10k & 0.47 (0.40--0.55) & 21 \\
50k / 15k & 0.49 (0.41--0.57) & 24 \\
\bottomrule
\end{tabular}
\end{table}
\clearpage

\noindent\textbf{Supplementary Table 4.} IRB institutions and protocol numbers (to be completed).

\end{document}
